%% GENERATED by notes/sysex-probes/sysex_examples_tex.py -- do not edit

\section{A parameter, written and read back}
The example is part~0's parameter \bytes{03}, a one-byte value, set to 100.

\begin{center}
\begin{tabular}{ll}
\toprule
& message \\
\midrule
write & \bytes{F0 50 2C 04 00 11 00 20 03 00 00 01 06 04 00 41 F7} \\
request & \bytes{F0 50 2B 04 00 11 00 20 03 00 00 01 00 4C F7} \\
reply & \bytes{F0 50 2C 04 00 11 00 20 03 00 00 01 06 04 00 41 F7} \\
part 7 & \bytes{F0 50 2C 04 00 11 00 27 03 00 00 01 06 04 00 3A F7} \\
three-byte & \bytes{F0 50 2C 04 00 11 00 20 00 00 00 03 01 02 00 00 04 00 00 45 F7} \\
\bottomrule
\end{tabular}
\end{center}

The value 100 is \bytes{64}, which the payload carries as the two bytes
\bytes{06 04}. The three bytes before it are the count, and the \bytes{00} after
it is the continuation flag. The checksum closes the message.

\section{A frame from a real instrument}
The bytes below open the first data message of a dump recorded from an
SX-WSA1R, a transfer of three categories in one session:

\begin{center}
\bytes{F0 50 2D 04 01 11 40 00 00 00 00 20 \ldots}
\end{center}

Reading it by the rules of chapter~4: family \bytes{2D}, model triple
\bytes{04 01 11} --- the rack's --- destination \bytes{40 00 00} and length \bytes{00 00 20},
which together name block~1 of \textsc{system, part \& midi}. The length decodes to 32 source
bytes and the frame carries 32 of them.

The payload decodes to \bytes{5A 5A 01 00 57 41 30 00}\ldots, the signature chapter~5 gives for
that block.

That block is the one chapter~\ref{ch:params} maps: a dump of this category can be
read against table~\ref{tbl:partparams} directly, which is the cheapest way to
confirm a parameter address without an instrument to hand.
